\section{引言}\label{sec:intro}

NPU 核内调度将一个算子级计算图中的数百至数万个微操作映射到有限的片上缓存层次上，这一过程延续了深度学习编译器中图级优化与多层 IR lowering 的基本路径~\citep{chen2018tvm,lattner2021mlir}。调度序必须满足数据依赖，并在片上缓存（scratchpad）的容量约束下管理中间值的驻留与换出。当同时驻留的数据量超出容量时，编译器需要将部分缓冲临时换出到片外内存并在后续重载，这类溢出既增加额外数据搬运，又可能阻塞执行流水线。

已有核内调度方法围绕关键路径、资源压力或峰值驻留量构造拓扑序优先级，将溢出视作容量不足后的局部后果。这一视角遗漏了一个结构性问题：当容量超限不可避免时，不同的合法调度序会在溢出窗口内暴露完全不同的缓冲构成，从而产生截然不同的换出代价。

具体而言，NPU 核内缓冲可自然区分为干净与脏两类：由片外读入且已有备份的干净缓冲在换出时无需写回，代价仅为重载；由计算产生的脏缓冲若后续仍被使用，则必须先写后读，代价恰为前者的两倍。二者在容量模型中占用相同的片上空间，却对应 $2\times$ 的溢出代价差异。若调度器只追求降低峰值或只在固定序上调整驱逐规则，就无法利用这一非对称结构。本文将这一通过改变合法拓扑序主动控制溢出窗口内 clean/dirty 构成的自由度称为溢出代价感知的活跃度塑形（spill-cost-aware liveness shaping）。

\begin{figure}[htbp]
\centering
\includegraphics[width=0.9\linewidth]{\ksFigure{e0_concept.pdf}}
\caption{溢出代价感知的活跃度塑形概念图。两个合法调度序可能具有相近的容量压力，但在溢出窗口内暴露给驱逐过程的 clean/dirty 构成不同。当 clean 缓冲在高压区保持常驻时，同等大小的换出可避免写回，从而降低额外搬运。}\label{fig:concept}
\end{figure}

本文的贡献如下：

\begin{enumerate}[leftmargin=2em]
  \item 溢出代价感知的活跃度塑形框架：通过调度序控制溢出窗口内的 clean/dirty 驻留构成，使高压区保留低成本可换出的缓冲储备。

  \item 溢出的理论刻画：溢出必然性定理（Theorem~\ref{thm:spill-inev}）给出线性时间可验证的单侧诊断条件，用于识别给定缓存容量下溢出是否为任何近最优调度都无法避免；Belady-margin 稳定性命题（Proposition~\ref{prop:belady-margin}）表明调度序而非驱逐策略是主要优化自由度；溢出面积近似定理（Theorem~\ref{thm:area-extra}）在有界缺席时长下建立 $A(S)$ 与最优溢出流量的双边常数界，解释代理目标在容量受限区域为何有效。

  \item 在公开 NPU 风格基准实例、合成分布、小图 oracle 和受控消融四个层面的系统实验，验证了本方法在缓存容量受限区域的有效性。
\end{enumerate}

实验表明，clean/dirty 无感的内存压力调度器会多付 $2.4\text{--}26\times$ 的额外溢出流量，且收益集中在定理判定为容量受限的区域。
